\documentclass[aspectratio=169]{beamer}
\usepackage[utf8]{inputenc}
\usepackage[T1]{fontenc}
\usetheme{solarized}
\usepackage{tabularx}
\usepackage[alf]{abntex2cite}

\title{Título da Apresentação}
\subtitle{Subtítulo}
\date{2023}
\author[Fulano]{Fulano de Tal}

% onde estão os códigos que aparecem nos slides
\lstset{inputpath=code}

\begin{document}

\maketitle

% Slide com o sumário
{
\def\insertframenumber{\emph{idx}}
\begin{frame}[t]\frametitle{\inserttitle}\framesubtitle{\insertsubtitle}
  % Use isso se tiver poucos tópicos
  % \tableofcontents[hidesubsections]

  % Ou use algo assim, se tiver muitos
  \begin{columns}[T]
    \begin{column}{.5\textwidth}
     \tableofcontents[hidesubsections,sections={1-4}]
    \end{column}
    \begin{column}{.5\textwidth}
      \tableofcontents[hidesubsections,sections={5-8}]
    \end{column}
  \end{columns}
  \addtocounter{framenumber}{-1}
\end{frame}
}

\section{Introdução}

\begin{frame}[c]\frametitle{\secname}\centering
  Esses são slides de exemplo de como usar o tema \inline*{beamerthemesolarized}.
  Esse tema é utiliza o modelo de cores Solarized disponível no seu próprio repositório do Github: \href{https://github.com/altercation/solarized}{altercation/solarized}.

  Os planos de fundo foram criados por mim no \href{https://inkscape.org/}{Inkscape} e exportados para PDF.
  Os slides devem ter proporção 16x9.
  Esse é um slide simples com um título, centralizado vertical e horizontalmente.
\end{frame}

\subsection{Títulos, subtítulos e alinhamento}

\begin{frame}[t,fragile]\frametitle{\secname}\framesubtitle{\subsecname}
  Para colocar título no slide, use \inline*{\textbackslash{}frametitle\{título\}}, para adicionar também um subtítulo, use \inline*{\textbackslash{}framesubtitle\{subtítulo\}}.

  Para alinhar o texto no topo, no centro ou em baixo use \inline*{[t]}, \inline*{[c]} ou \inline*{[b]}.

  Para centralizar o texto use \inline*{\textbackslash{}centering} e para alinhar à direita use \inline*{\textbackslash{}raggedleft} que quer dizer que ele vai ficar picotado à esquerda.
  O padrão é alinhar à esquerda e ficar picotado à direita.

  Esse slide foi criado com um código mais ou menos assim:
  \begin{snippet}[style=text,literate={\#}{{}}{0}]
\begin{frame}[t]%
\frametitle{\secname}\framesubtitle{\subsecname}
  Para ...
\#end{frame}
  \end{snippet}
\end{frame}

\section{Seções}

\subsection{Também subseções e subsubseções}
\begin{frame}[t]\frametitle{\secname}\framesubtitle{\subsecname}
  Para iniciar uma nova seção dos slides use o comando \inline*{\textbackslash{}section\{seção\}}.
  Esse comando criará um slide de título de seção como o slide anterior a esse.
  Para criar uma subseção utilize o comando \inline*{\textbackslash{}subsection\{seção\}}.
  Os nomes da seção e subseção atuais ficam disponíveis, respectivamente, nas macros \inline*{\textbackslash{}secname} e \inline*{\textbackslash{}subsecname}, podendo serem usados para colocar título e subtítulo nos slides.

  \subsubsection{Subsubseções}

  Há ainda um terceiro nível de seção que pode ser criado com o comando \inline*{\textbackslash{}subsubsection\{Subsubseção\}}. Esse comando cria um divisão dentro do próprio slide.
\end{frame}

\section{Estilos Escolhidos}

\subsection{Fonte e fundo}

\begin{frame}[t]\frametitle{\secname}\framesubtitle{\subsecname}
  A fonte tipográfica escolhida para esse modelo foi a Fira Sans.
  Além disso, os ``fundos'' dos slides foram criados em SVG no software Inkscape usando as cores do tema Solarized e posteriormente convertidos para PDF.

  Esses fundos são apenas o rodapé dos slides.
  Na verdade o fundo está posicionado \textbf{sobre} o restante do conteúdo.
  Isso para evitar cobrir cobrir os detalhes do fundo. Além disso,
  interromper o conteúdo dos slides na altura do topo da imagem do fundo reduziria muito a quantidade de conteúdo que poderia ser apresentada nos slides.
  Por isso, optou-se por permitir que o conteúdo flua por trás da imagem de ``fundo''.
  
  Isso fica bem aparente quando usamos caixas coloridas, imagens ou tabelas que descem até embaixo, como essa aqui:

  \color{red}\rule{\textwidth}{0.5cm}
\end{frame}

\subsection{Slides de título e seção}

\begin{frame}[t]\frametitle{\secname}\framesubtitle{\subsecname}
  Os slides de título e de seção também seguem um modelo semelhante.
  Nesse caso, os detalhes à direita do rodapé do slide são maiores e um pouco diferentes.

  Nos slides de título e de seção há ainda o título e subtítulo da apresentação centralizados na parte de cima.
\end{frame}

\subsection{Cores}

\setbeamercolor{white}{bg=white}
\begin{frame}[t]\frametitle{\secname}\framesubtitle{\subsecname}
  As seguintes cores estão disponíveis com seus próprios nomes.
  As cores da linha de cima são recomendadas para fundo e para conteúdo padrão, as da linha de baixo para realce.\\[1em]
  \begin{beamercolorbox}[center]{white}
    ~\\[1em]\scriptsize
  \begin{tabular}{cccccccc}
    \color{base03}\rule{1.2cm}{1.2cm} & \color{base02}\rule{1.2cm}{1.2cm} & \color{base01}\rule{1.2cm}{1.2cm} & \color{base00}\rule{1.2cm}{1.2cm} & \color{base0}\rule{1.2cm}{1.2cm}  & \color{base1}\rule{1.2cm}{1.2cm}  &\color{base2}\rule{1.2cm}{1.2cm}  & \color{base3}\rule{1.2cm}{1.2cm} \\
    base03 & base02 & base01 & base00 & base0 & base1 & base2 & base3 \\[1em]
    \color{yellow}\rule{1.2cm}{1.2cm} & \color{orange}\rule{1.2cm}{1.2cm} & \color{red}\rule{1.2cm}{1.2cm} & \color{magenta}\rule{1.2cm}{1.2cm} & \color{violet}\rule{1.2cm}{1.2cm}  & \color{blue}\rule{1.2cm}{1.2cm}  &\color{cyan}\rule{1.2cm}{1.2cm}  & \color{green}\rule{1.2cm}{1.2cm} \\
    yellow & orange & red & magenta & violet & blue & cyan & green \\[1em]
  \end{tabular}
  \end{beamercolorbox}
\end{frame}

\section{Texto}

\subsection{Ênfase}

\begin{frame}[t,fragile]\frametitle{\secname}\framesubtitle{\subsecname}
  Para dar ênfase, o texto pode ser escrito em \textbf{negrito}, \emph{itálico}, ou em \alert{vermelho}, como o usual.

  As cores de realce de texto estão disponíveis como comandos em português:
  \amarelo{amarelo} \laranja{laranja} \vermelho{vermelho} \magenta{magenta} \violeta{violeta} \azul{azul} \ciano{ciano} \verde{verde}.

  Outras opções de ênfase também são alcançáveis através de combinações de comandos simples.
  O texto pode ser escrito em \cianito{``cianito''}, por exemplo, porque esse comando já foi definido como:
\begin{snippet}[style=text]
\newcommand{\cianito}[1]{%
  \ciano{\textbf{#1}}%
}
\end{snippet}

\end{frame}

\subsection{Colunas}
\begin{frame}[t]\frametitle{\secname}\framesubtitle{\subsecname}
  \begin{columns}[T]
    \begin{column}{.5\textwidth}
      \subsubsection{Coluna da esquerda}

      O texto também pode ser quebrado em colunas.
      O ideal seria evitar quebrar em muitas colunas.
    \end{column}
    \begin{column}{.5\textwidth}
      \subsubsection{Coluna da direita}

      Para blocos de texto, duas colunas é o suficiente.
      Passar disso fará com que as colunas fiquem estreitas.
    \end{column}
  \end{columns}
  \subsubsection{Para listas podemos usar mais}
  \begin{columns}[T]
    \begin{column}{.25\textwidth}
      \begin{itemize}
        \item Item A
        \item Item B
        \item Item C
      \end{itemize}
    \end{column}
    \begin{column}{.25\textwidth}
      \begin{itemize}
        \item Item D
        \item Item E
        \item Item F
      \end{itemize}
    \end{column}
    \begin{column}{.25\textwidth}
      \begin{itemize}
        \item Item G
        \item Item H
        \item Item I
      \end{itemize}
    \end{column}
    \begin{column}{.25\textwidth}
      \begin{itemize}
        \item Item J
        \item Item K
        \item Item L
      \end{itemize}
    \end{column}
  \end{columns}
\end{frame}

\subsection{Citações longas}
\begin{frame}[t]\frametitle{\secname}\framesubtitle{\subsecname}
  Citações longas podem usar o ambiente \inline*{quotation} que vai formatar e colocar aspas tipográficas na volta para efeito estético:

  \begin{quotation}
    Pra falar bem francamente: enquanto não havia máquinas, programar não era nenhum problema.
  \end{quotation}
  --- Edsger Dijkstra\raggedleft
\end{frame}

\subsection{Blocos}

\begin{frame}[t]\frametitle{\secname}\framesubtitle{\subsecname}
  \begin{block}{Título do bloco}
    Texto do bloco.
  \end{block}
  \begin{exampleblock}{Título do exemplo}
    Texto do exemplo.
  \end{exampleblock}
  \begin{alertblock}{Título do alerta}
    Texto do alerta.
  \end{alertblock}
  {
    \setbeamercolor{block title}{bg=yellow}
    \setbeamercolor{block body}{fg=yellow}
    \begin{block}{Personalizado}
      Exemplo de bloco com cor personalizada.
    \end{block}
  }
\end{frame}

\subsection{Tabelas}

\begin{frame}[t]\frametitle{\secname}\framesubtitle{\subsecname}
  Tabelas podem ser inseridas em qualquer slide.
  Essa é uma tabela bem completa com legenda e link para a fonte dos dados:
  \begin{table}
    \caption{\label{tab:pib}PIB dos estados da região sul}
    \begin{tabular}{l r r r r}
      \toprule
      \cianito{UF} & \multicolumn{1}{c}{\cianito{PIB}} & \cianito{\% NACIONAL} & \cianito{\% REGIONAL} & \cianito{PER CAPTA}\\
      \midrule
      PR & 136 681 933 mil & 6,4 & 34,2 & 13 158,00 \\
      SC & 93 193 324 mil & 4,0 & 21,5 & 15 638,00 \\
      RS & 156 883 171 mil & 8,2 & 44,3 & 14 310,00 \\
      \bottomrule
    \end{tabular}\\[1ex]
    \makebox[12cm][r]{\small Valores em reais. FONTE: \href{https://pt.wikipedia.org/wiki/Regi\%C3\%A3o_Sul_do_Brasil}{Wikipedia}.}
  \end{table}
\end{frame}

\begin{frame}[t]\frametitle{\secname}\framesubtitle{\subsecname}
  A mesma tabela poderia ser apresentada de maneira mais ``limpa'', facilitando que a audiência foque sua atenção nos dados:\footnote[frame]{PIB dos estados da região sul. Valores em reais. FONTE: \href{https://pt.wikipedia.org/wiki/Regi\%C3\%A3o_Sul_do_Brasil}{Wikipedia}}
  \begin{table}
    \begin{tabular}{l r r r r}
      \toprule
      \cianito{UF} & \multicolumn{1}{c}{\cianito{PIB}} & \cianito{\% NACIONAL} & \cianito{\% REGIONAL} & \cianito{PER CAPTA}\\
      \midrule
      PR & 136 681 933 mil & 6,4 & 34,2 & 13 158,00 \\
      SC & 93 193 324 mil & 4,0 & 21,5 & 15 638,00 \\
      RS & 156 883 171 mil & 8,2 & 44,3 & 14 310,00 \\
      \bottomrule
    \end{tabular}
  \end{table}
\end{frame}

\section{Imagens}

\begin{frame}[t]\frametitle{\secname}
  Vários tipos de imagens são suportados.
  Os seguintes são recomendados:
  \def\arraystretch{1.5}
  \begin{table}
    \begin{tabular}{l c}
      \toprule
      \cianito{CONTEÚDO} & \cianito{TIPOS RECOMENDADOS}\\
      \midrule
      Linhas, diagramas, gráficos, texto, etc. & PDF ou EPS\\
      Figuras com cores sólidas, capturas de telas, etc. & PNG ou TIFF\\
      Fotografias, imagens com muitos detalhes, etc. & JPG\\
      \bottomrule
    \end{tabular}
  \end{table}
\end{frame}

\subsection{Legenda ao lado}

\begin{frame}[t]\frametitle{\secname}\framesubtitle{\subsecname}
  \begin{columns}[T]
    \begin{column}{9.75cm}
      \includegraphics[width=\columnwidth]{images/pdf/complexidade.pdf}%
    \end{column}
    \begin{column}{4.25cm}
      \it\cianito{complexidade temporal}\\
      imagem gerada com o \href{http://www.gnuplot.info/}{Gnuplot}, terminal cairolatex.
      Essa é uma imagem na proporção 6x4, com 9,75x6,50cm. Colocar a legenda ao lado permite uma legenda maior.
    \end{column}
  \end{columns}
\end{frame}

\begin{frame}[t]\frametitle{\secname}\framesubtitle{\subsecname}
  \begin{columns}[T]
    \begin{column}{9.75cm}
      \includegraphics[width=\columnwidth]{images/pdf/switch.pdf}%
      \par\vspace*{-1.8\baselineskip}\hfill{}\colorbox{base2}{\footnotesize\color{base02}Créditos: \href{https://github.com/roger-willian}{Autor}}
    \end{column}
    \begin{column}{4.25cm}
      \it\cianito{switch}\\
      imagem SVG (\firalining 75 KB) gerada com o \href{https://inkscape.org/}{Inkscape} e exportada para PDF (\firalining 335 KB).
      É possível utilizar \LaTeX{} dentro do SVG para formatar equações, se for necessário.
      O PDF mantém o formato vetorial, amplie e observe as letras.
    \end{column}
  \end{columns}
\end{frame}

\subsection{Vetorial vs. bitmap}

\begin{frame}[t]\frametitle{\secname}\framesubtitle{\subsecname}
  \begin{columns}[T]
    \begin{column}{9.75cm}
      \includegraphics[width=\columnwidth]{images/bitmap/switch.png}%
      \par\vspace*{-1.8\baselineskip}\hfill{}\colorbox{base2}{\footnotesize\color{base02}Créditos: \href{https://github.com/roger-willian}{Autor}}
    \end{column}
    \begin{column}{4.25cm}
      \it\cianito{switch}\\
      mesma imagem exportada para o formato PNG (\firalining 277 KB).
      Esse \textbf{não} é um formato vetorial, amplie e observe como agora as letras ficaram pixeladas.
    \end{column}
  \end{columns}
\end{frame}

\subsection{Legenda abaixo}

\begin{frame}[t]\frametitle{\secname}\framesubtitle{\subsecname}
  \begin{columns}[T]
    \begin{column}{\textwidth}
      \includegraphics[width=\columnwidth]{images/pdf/etapas.pdf}%
      \\\small
      \it\cianito{processo} --- 7x3 (14x6cm), gerada com \href{https://app.diagrams.net/}{diagrams.net}.
    \end{column}
  \end{columns}
\end{frame}

\section{Código Fonte}

\subsection{Inline}

\lstset{language=C}
\begin{frame}[t]\frametitle{\secname}\framesubtitle{\subsecname}
  Para colocar um pequeno pedaço de código com fonte monoespaçada e realce de cores no meio de uma linha de texto normal, use o comando \inline*{\textbackslash{}inline\{\}}.
  Por exemplo:

  ``A função \inline{printf()} da biblioteca padrão da linguagem C serve para imprimir uma string formatada na tela.''

  Já para colocar um texto com fonte monoespaçada, mas sem realce de cores, use o comando \inline*{\textbackslash{}inline*\{\}} com asterisco. Por exemplo:

  ``O arquivo de cabeçalho \inline*{stdio.h} contém as declarações das funções de entrada e saída.''
\end{frame}

\subsection{Tabelas}

\begin{frame}[t]\frametitle{\secname}\framesubtitle{\subsecname}
  Para colocar fonte monoespaçada em uma coluna de uma tabela, são fornecidas duas colunas especiais \inline*{I} e \inline*{i} que usam os comandos \inline*{\textbackslash{}inline\{\}} e \inline*{\textbackslash{}inline*\{\}} automaticamente:

  ~
  \def\arraystretch{1.5}
  \centering
  \begin{tabular}{I i l}
    \toprule
    \multicolumn{1}{l}{\cianito{Função}} & \multicolumn{1}{l}{\cianito{Biblioteca}} & \cianito{Descrição} \\
    \midrule
    printf() & <stdio.h> & Imprime string formatada \\
    getchar() & <stdio.h> & Lê um caractere do teclado \\
    memcpy() & <string.h> & Copia bloco de memória \\
    malloc() & <stdlib.h> & Aloca bloco de memória \\
    \bottomrule
  \end{tabular}
\end{frame}

\subsection{Snippet}

\begin{frame}[t,fragile]\frametitle{\secname}\framesubtitle{\subsecname}
  Para colocar um trecho maior de código, ocupando uma ou mais linhas, use o ambiente \inline*{snippet}. 
  No entanto, é preciso passar a opção \inline*{fragile} para o \inline*{frame}. Por exemplo:

  \begin{snippet}
int fatorial(int n) {
  if (n < 0) return -1;

  int fat_n = 1;
  for(int i = 1; i <= n; i++) {
    fat_n *= i;
  }
  return fat_n;
}
  \end{snippet}
\end{frame}

\begin{frame}[t,fragile]\frametitle{\secname}\framesubtitle{\subsecname}
  Tanto o comando \inline*{\textbackslash{}inline\{\}} quanto o ambiente snippet usam o pacote \inline*{listings}\footnote[frame]{Documentação na \href{http://mirrors.ctan.org/macros/latex/contrib/listings/listings.pdf}{CTAN}} para criar o realce de sintaxe.
  Porém o ambiente \inline*{snippet} aceita um parâmetro opcional que são as opções para o ambiente \inline*{lstlistings}.
  Por exemplo, o trecho a seguir usa uma fonte gigante:

\begin{snippet}[basicstyle=\Huge\ttfamily\color{codefg}]
int maior(int x, int y) {
  return x > y ? x : y;
}
\end{snippet}
\end{frame}

\begin{frame}[t,fragile]\frametitle{\secname}\framesubtitle{\subsecname}
  Usando essas opções também é possível mostrar um trecho de texto em fonte monoespaçada sem realce de cores:
\begin{snippet}[style=text]
LISTA DO MERCADO
Arroz       5 kg
Feijão      2 kg
Leite       6 cx
Pão         1 pct
\end{snippet}
\end{frame}

\subsection{Falso editor}

\begin{frame}[t]\frametitle{\secname}\framesubtitle{\subsecname}
  Para inserir um arquivo como se fosse dentro de um editor, use o comando \inline*{\textbackslash{}arquivo\{\}} que usa os pacotes \inline*{listings} para o realce e \inline*{tcolorbox} para criar a ilusão da janela.\footnote[frame]{Documentação na \href{http://mirrors.ctan.org/macros/latex/contrib/tcolorbox/tcolorbox.pdf}{CTAN}} Por exemplo:
  \arquivo{ola-mundo.c}

  Importante: o nome do arquivo não pode ter caracteres especiais do \LaTeX{}  como o caractere de sublinhado \inline*{\_}.
\end{frame}

\begin{frame}[t]\frametitle{\secname}\framesubtitle{\subsecname}
  Esse comando pode receber ainda dois parâmetros opcionais:
  \begin{itemize}
    \item antes do nome do arquivo, as opções para o ambiente \inline*{lstlistings};
    \item depois do nome do arquivo, as opções para o ambiente \inline*{tcolorbox};
  \end{itemize}

  Por exemplo, mostrar apenas algumas linhas e trocar o título da janela:
\arquivo[firstline=3,firstnumber=3]{ola-mundo.c}[title={Trecho do mesmo arquivo}]
\end{frame}

\subsection{Slide todo}

\begin{frame}[t]\frametitle{\secname}\framesubtitle{\subsecname}
  Um slide normal pode ser pequeno para mostrar um trecho um pouco maior de um arquivo.
  O comando \inline*{\textbackslash{}arquivo\{\}} consegue mostrar no máximo umas 15 linhas.

  Por outro lado, é possível utilizar um slide inteiro para mostrar o arquivo (veja o próximo slide).
  Basta utilizar o comando \inline*{\textbackslash{}framearquivo\{\}}.
  Esse comando consegue mostrar 21 linhas de um arquivo.

  O comando \inline*{\textbackslash{}framearquivo\{\}} tem basicamente o mesmo formato que o comando \inline*{\textbackslash{}arquivo\{\}}, incluindo os dois parâmetros opcionais.
  Contudo, esse comando deve ser usado fora de qualquer slide porque ele cria um novo slide por si só.
\end{frame}

\framearquivo{fatorial.c}

\subsection{Falso terminal}

\begin{frame}[t]\frametitle{\secname}\framesubtitle{\subsecname}
  Também é possível simular um terminal, caso seja preciso falar sobre algum comando:

  ~
  \def\prompt{\color{codekeyword}usuario@maquina:{\textasciitilde}\$\color{codefg}~}
  \begin{terminal}[title={Bash}]
  \prompt ls
  Desktop\\Documents\\Downloads\\Music\\Pictures\\Public\\Videos\\
  \prompt \rule{1ex}{1em}
  \end{terminal}
\end{frame}

\subsection{Tema light vs. dark}

\begin{frame}[t]\frametitle{\secname}\framesubtitle{\subsecname}
  Dois esquemas de cores estão disponíveis para o realce de sintaxe do código, ambos baseados no esquema solarized\footnote[frame]{Disponível no \href{https://ethanschoonover.com/solarized/}{Github}}.
  O esquema de cores pode ser trocado usando os comandos \inline*{\textbackslash{}codelight} (padrão) ou \inline*{\textbackslash{}codedark\{\}}:

  ~
  \begin{columns}[T]
    \begin{column}{.33\textwidth}
      \begin{tabularx}{\columnwidth}{c X}
        \toprule
        \multicolumn{2}{c}{\cianito{Ambos}} \\
        \midrule
        \color{codekeyword}   \rule[-0.5ex]{0.5cm}{0.5cm} & codekeyword   \\
        \color{codevar}       \rule[-0.5ex]{0.5cm}{0.5cm} & codevar       \\
        \color{codestring}    \rule[-0.5ex]{0.5cm}{0.5cm} & codestring    \\
        \color{codedirective} \rule[-0.5ex]{0.5cm}{0.5cm} & codedirective \\
        \color{codecomment}   \rule[-0.5ex]{0.5cm}{0.5cm} & codecomment   \\
        \bottomrule
      \end{tabularx}
    \end{column}
    \begin{column}{.33\textwidth}
      \begin{tabularx}{\columnwidth}{c X}
        \toprule
        \multicolumn{2}{c}{\cianito{Light}} \\
        \midrule
        \color{codebg}    \rule[-0.5ex]{0.5cm}{0.5cm} & codebg    \\
        \color{codefg}    \rule[-0.5ex]{0.5cm}{0.5cm} & codefg    \\
        \color{codelines} \rule[-0.5ex]{0.5cm}{0.5cm} & codelines \\
        \color{codetitle} \rule[-0.5ex]{0.5cm}{0.5cm} & codetitle \\
        \color{codetext}  \rule[-0.5ex]{0.5cm}{0.5cm} & codetext  \\
        \bottomrule
      \end{tabularx}
    \end{column}
    \begin{column}{.33\textwidth}\codedark
      \begin{tabularx}{\columnwidth}{c X}
        \toprule
        \multicolumn{2}{c}{\cianito{Dark}} \\
        \midrule
        \color{codebg}    \rule[-0.5ex]{0.5cm}{0.5cm} & codebg    \\
        \color{codefg}    \rule[-0.5ex]{0.5cm}{0.5cm} & codefg    \\
        \color{codelines} \rule[-0.5ex]{0.5cm}{0.5cm} & codelines \\
        \color{codetitle} \rule[-0.5ex]{0.5cm}{0.5cm} & codetitle \\
        \color{codetext}  \rule[-0.5ex]{0.5cm}{0.5cm} & codetext  \\
        \bottomrule
      \end{tabularx}
    \end{column}
  \end{columns}
\end{frame}

\begin{frame}[t]\frametitle{\secname}\framesubtitle{\subsecname}
  \centering\cianito{Exemplo: Inline}

  ~
  \begin{columns}[T]
    \begin{column}{.5\textwidth}
      \codelight
      O esquema claro é melhor para a apresentação em uma tela projetada.

      Esse é um texto de exemplo com realce de sintaxe:
      
      \inline{int x = 10;}

      Já esse outro exemplo não tem realce:
      
      \inline*{meuarquivo.txt}
    \end{column}
    \begin{column}{.5\textwidth}
      \codedark
      O esquema escuro é melhor para ler a apresentação diretamente no PC.

      Esse é um texto de exemplo com realce de sintaxe:
      
      \inline{int x = 10;}

      Já esse outro exemplo não tem realce:
      
      \inline*{meuarquivo.txt}
    \end{column}
  \end{columns}
\end{frame}

\begin{frame}[t,fragile]\frametitle{\secname}\framesubtitle{\subsecname}
  \centering\cianito{Exemplo: Snippet}

  ~
  \begin{columns}[T]
    \begin{column}{.5\textwidth}
      \codelight
\begin{snippet}
int maior(int a, int b) {
  return a > b ? a : b;
}
\end{snippet}
    \end{column}
    \begin{column}{.5\textwidth}
      \codedark
\begin{snippet}
int maior(int a, int b) {
  return a > b ? a : b;
}
\end{snippet}
    \end{column}
  \end{columns}
\end{frame}

\begin{frame}[t]\frametitle{\secname}\framesubtitle{\subsecname}
  \centering\cianito{Exemplo: Arquivos}

  ~
  \begin{columns}[T]
    \begin{column}{.5\textwidth}
      \codelight
      \arquivo[firstline=3,lastline=11,firstnumber=3]{fatorial.c}[]
    \end{column}
    \begin{column}{.5\textwidth}
      \codedark
      \arquivo[firstline=3,lastline=11,firstnumber=3]{fatorial.c}[]
    \end{column}
  \end{columns}
\end{frame}

\framearquivo{fatorial.c}

{\codedark\framearquivo{fatorial.c}}

\section{Listas}

\subsection{Definições}

\begin{frame}\frametitle{\secname}\framesubtitle{\subsecname}\label{slide:frutas}
  \subsubsection{Às vezes é preciso uma lista de definições:}
  \begin{columns}[T]
    \begin{column}{0.48\textwidth}
      \begin{description}[abacaxi]
        \item[Abacate] fruto do abacateiro.
        \item[Abacaxi] espécie de ananás.
        \item[Açaí] fruta do açaizeiro.
        \item[Acerola] fruto semelhante à cereja.
        \item[Amora] fruto da amoreira e de algumas espécies de silvas.
        \item[Araçá] fruto de uma árvore myrtácea da América.
      \end{description}
    \end{column}\hfill{}\color{orange}\vline\hfill{}
    \begin{column}{0.48\textwidth}
      \begin{description}[Carambola]
        \item[Araticum] designação de várias árvores do Brasil; fruto dessa árvore.
        \item[Banana] fruto da bananeira.
        \item[Cáqui] fruto amarelo e comestível.
        \item[Carambola] fruto do caramboleiro.
      \end{description}
    \end{column}
  \end{columns}
\end{frame}

\subsection{Listas}

\begin{frame}[t]\frametitle{\secname}\framesubtitle{\subsecname}
  As listas são configuradas para ficar mais próximo do restante do tema:

  ~
  \begin{columns}[T]
    \begin{column}{.5\textwidth}
      Lista não ordenada:
      \begin{itemize}
        \item Cachorro-quente;
        \item Refrigerante;
        \item Fritas.
      \end{itemize}
    \end{column}
    \begin{column}{.5\textwidth}
      Lista ordenada:
      \begin{enumerate}
        \item Jantar;
        \item Escovar os dentes;
        \item Dormir.
      \end{enumerate}
    \end{column}
  \end{columns}
\end{frame}

\subsection{Lista longa}

\begin{frame}[t]\frametitle{\secname}\framesubtitle{\subsecname}
  Uma lista longa continua a sua numeração automaticamente:
  \begin{columns}[T]
    \begin{column}{.5\textwidth}
      \begin{listalonga*}
        \item O quê?
        \item Quando?
        \item Por que?
      \end{listalonga*}
    \end{column}
    \begin{column}{.5\textwidth}
      \begin{listalonga}
        \item Quem?
        \item Como?
        \item Onde?
      \end{listalonga}
    \end{column}
  \end{columns}
  ~\\

  Para iniciar uma nova lista longa, utilize o ambiente \inline*{listalonga*}:
  \begin{columns}[T]
    \begin{column}{.5\textwidth}
      \begin{listalonga*}
        \item Quedêle?
        \item Quié?
        \item Cumé?
      \end{listalonga*}
    \end{column}
    \begin{column}{.5\textwidth}
      ~
    \end{column}
  \end{columns}
\end{frame}

\section{Outros}

\subsection{Expressões matemáticas}

\newcommand{\pulleqnum}[1][1cm]{\stepcounter{equation}\tag*{(\theequation)\hspace{#1}}}
\begin{frame}[t]\frametitle{\secname}\framesubtitle{\subsecname}
  É possível usar equações no meio de um slide, o volume $v$ de uma esfera é dado por:
  \begin{align}\label{eq:volume}
    v &= \int_0^r \int_0^{2\pi} \int_0^\pi \rho^2 \sin(\theta)\, d\theta\, d\varphi\, d\rho \\[1em]
      &= \frac{4}{3}\pi r^3
  \end{align}

  Aqui, o determinante de uma matriz singular:
  \begin{equation}
    \det \begin{pmatrix}
      1 & 2 & 3 \\
      4 & 5 & 6 \\
      7 & 8 & 9
    \end{pmatrix} = 0\pulleqnum
  \end{equation}
\end{frame}

\subsection{Citações e referências}

\begin{frame}[t]\frametitle{\secname}\framesubtitle{\subsecname}%
  \subsubsection{Da literatura}
  Quando necessário pode-se colocar uma citação indicando a referência bibliográfica original de onde uma ideia foi retirada.
  Por exemplo, há uma norma brasileira que rege as citações e referências bibliográficas \cite{NBR6023:2000}.

  \subsubsection{De outro slide}
  É possível fazer referência a outras partes dessa mesma apresentação de slides também.
  A expressão do volume em \eqref{eq:volume} é ``facilmente'' derivada, mas as definições de frutas apresentadas anteriormente no Slide~\slideref{slide:frutas} foram\\ retiradas de um dicionário \cite{figueiredo1913novo}.
\end{frame}

\subsection{Agradecimento}

\begin{frame}[t]\frametitle{\secname}\framesubtitle{\subsecname}
  Também é fornecido um comando específico para produzir um slide final de agradecimento.
  Esse comando é o \inline*{\textbackslash{}agradecimento\{\}} que recebe como parâmetro o agradecimento.
  Ele gera um slide semelhante ao slide de título, mas com o agradecimento em baixo.

  Por exemplo, o próximo slide foi produzido com o comando

  \inline*{\textbackslash{}agradecimento\{Obrigado pela oportunidade!\}}
  
\end{frame}

\agradecimento{Obrigado pela oportunidade!}

\section*{Referências}
\setbeamertemplate{bibliography item}{\rule{0pt}{1.5em}}
\setbeamertemplate{footline}{%
    \let\velhoframenumber\insertframenumber
    \def\insertframenumber{R-\velhoframenumber}
    \vspace{-1.5cm}\hfill\usebeamertemplate***{page number in head/foot}\hspace*{-5mm}
    \vspace{3mm}
}
\begin{frame}[t]\frametitle{\secname}
    \footnotesize
    \bibliography{referencias.bib}
\end{frame}

\end{document}